\documentclass[11pt]{article}
% ===========================================================================
%  Shared preamble for all MPM2D worksheets — input right after \documentclass.
%  (Files beginning with "_" are skipped by mpm2d-worksheet-publish.mjs.)
% ===========================================================================
\usepackage[letterpaper,top=2.2cm,bottom=1.9cm,left=1.6cm,right=1.6cm,headheight=28pt,footskip=24pt]{geometry}
\usepackage{amsmath,amssymb}
\usepackage[T1]{fontenc}
\usepackage[utf8]{inputenc}
\usepackage{lmodern}
\usepackage{xcolor}
\usepackage[most]{tcolorbox}
\usepackage{enumitem}
\usepackage{multicol}
\usepackage{fancyhdr}
\usepackage{tikz}
\usetikzlibrary{calc}
\usepackage{pgfplots}
\pgfplotsset{compat=1.18}

\definecolor{exblue}{HTML}{4A90E2}
\definecolor{exbg}{HTML}{E6F3FF}
\definecolor{qorange}{HTML}{E69138}
\definecolor{qbg}{HTML}{FFF7CC}
\definecolor{keygray}{HTML}{475569}
\definecolor{keybg}{HTML}{F1F5F9}
\definecolor{iagreen}{HTML}{14653B}
\definecolor{titleblue}{HTML}{1D4ED8}
% Rotating palette so each worked example gets its own colour.
\definecolor{ex0}{HTML}{2563A0}\definecolor{ex1}{HTML}{3B7D3B}\definecolor{ex2}{HTML}{8A5A00}
\definecolor{ex3}{HTML}{A3327A}\definecolor{ex4}{HTML}{6D28A3}\definecolor{ex5}{HTML}{B08900}
\definecolor{ex6}{HTML}{0E7490}\definecolor{ex7}{HTML}{9A3412}\definecolor{ex8}{HTML}{15803D}

\pagestyle{fancy}
\fancyhf{}
\fancyhead[L]{\footnotesize NAME: \underline{\hspace{3.4cm}}}
\fancyhead[C]{\textbf{Integration Academy}\\[-2pt]{\scriptsize Math Simplified \textperiodcentered{} Success Amplified}}
\fancyhead[R]{\footnotesize MPM2D}
\fancyfoot[L]{\scriptsize dr.merie@IntegrationAcademy.ca}
\fancyfoot[C]{\scriptsize\thepage}
\fancyfoot[R]{\scriptsize IntegrationAcademy.ca}
\renewcommand{\headrulewidth}{1.2pt}
\renewcommand{\footrulewidth}{0.4pt}
\setlength{\parindent}{0pt}
\setlength{\parskip}{4pt}

% Examples are NOT breakable, so each example + its solution stays on one page.
\newtcolorbox{example}[2]{colframe=#1,colback=#1!7,coltitle=white,colbacktitle=#1,fonttitle=\bfseries,title={#2},arc=3pt,boxrule=1pt}
\newtcolorbox{qbox}[1]{colframe=qorange,colback=qbg,coltitle=white,colbacktitle=qorange,fonttitle=\bfseries,title={#1},arc=3pt,breakable,boxrule=1pt}
\newtcolorbox{keybox}{colframe=keygray,colback=keybg,coltitle=white,colbacktitle=keygray,fonttitle=\bfseries,title={$\checkmark$\ Answer Key},arc=3pt,breakable,boxrule=1pt}
\newtcolorbox{recapbox}{colframe=keygray,colback=keybg,coltitle=white,colbacktitle=keygray,fonttitle=\bfseries,title={Question Recap},arc=3pt,breakable,boxrule=1pt}
\newtcolorbox{ideabox}{colframe=titleblue,colback=titleblue!6,coltitle=white,colbacktitle=titleblue,fonttitle=\bfseries,title={The Big Ideas},arc=3pt,breakable,boxrule=1.2pt}
\newtcolorbox{learnbox}{colframe=iagreen,colback=iagreen!5,coltitle=white,colbacktitle=iagreen,fonttitle=\bfseries,title={Concept Essentials},arc=3pt,breakable,boxrule=1.2pt}
\newcommand{\lh}[1]{\par\smallskip{\bfseries\color{iagreen}#1}\par\nobreak\smallskip}

\newcommand{\soln}{\par\smallskip\textit{\textbf{Solution.}}\ }
\newcommand{\workspace}[1]{\par\vspace{3pt}\fcolorbox{black!30}{white}{\begin{minipage}[t][#1][t]{\dimexpr\linewidth-2\fboxsep\relax}\ \end{minipage}}\par}
\newcommand{\sech}[1]{\par\vspace{8pt}{\Large\bfseries\color{iagreen}#1}\par\nobreak\vspace{2pt}{\color{black!20}\hrule height1pt}\vspace{6pt}}

% A compact coordinate plot sized for an example box: \eplot{xmin}{xmax}{ymin}{ymax}{body}
\newcommand{\eplot}[5]{\begin{center}\begin{tikzpicture}
\begin{axis}[width=6.8cm,height=4.4cm,axis lines=middle,xlabel={\tiny$x$},ylabel={\tiny$y$},
grid=both,grid style={gray!16},xmin=#1,xmax=#2,ymin=#3,ymax=#4,
xtick distance=2,ytick distance=2,tick label style={font=\tiny},axis line style={-},samples=120]
#5
\end{axis}\end{tikzpicture}\end{center}}

% A sine-curve plot (degrees on x, 0..360): \splot{ymin}{ymax}{body}
\newcommand{\splot}[3]{\begin{center}\begin{tikzpicture}
\begin{axis}[width=8.6cm,height=4.6cm,axis lines=middle,xlabel={\tiny$x\,(^\circ)$},ylabel={\tiny$y$},
grid=both,grid style={gray!16},xmin=0,xmax=360,ymin=#1,ymax=#2,
xtick distance=90,ytick distance=1,tick label style={font=\tiny},axis line style={-},samples=180]
#3
\end{axis}\end{tikzpicture}\end{center}}

% A small coordinate plot: \plot{xmin}{xmax}{ymin}{ymax}{body}
\newcommand{\plot}[5]{\begin{center}\begin{tikzpicture}
\begin{axis}[width=8.4cm,height=6cm,axis lines=middle,xlabel={\scriptsize$x$},ylabel={\scriptsize$y$},
grid=both,grid style={gray!16},xmin=#1,xmax=#2,ymin=#3,ymax=#4,
xtick distance=2,ytick distance=2,tick label style={font=\tiny},axis line style={-}]
#5
\end{axis}\end{tikzpicture}\end{center}}

% Right triangle (angle theta at bottom-left, right angle at bottom-right):
%   \rtri{drawBase}{drawHeight}{oppLabel}{adjLabel}{hypLabel}
\newcommand{\rtri}[5]{\begin{center}\begin{tikzpicture}[scale=0.85]
\coordinate (A) at (0,0);\coordinate (B) at (#1,0);\coordinate (C) at (#1,#2);
\draw[thick,exblue] (A)--(B)--(C)--cycle;
\draw[black] ($(B)+(-0.32,0)$)--($(B)+(-0.32,0.32)$)--($(B)+(0,0.32)$);
\node[below] at ($(A)!0.5!(B)$){#4};
\node[right] at ($(B)!0.5!(C)$){#3};
\node[above left] at ($(A)!0.5!(C)$){#5};
\node at (0.7,0.22){$\theta$};
\end{tikzpicture}\end{center}}

% General (oblique) triangle with vertex labels and opposite-side labels:
%   \gtri{Alabel}{Blabel}{Clabel}{aSide}{bSide}{cSide}
\newcommand{\gtri}[6]{\begin{center}\begin{tikzpicture}[scale=0.85]
\coordinate (A) at (0,0);\coordinate (B) at (5,0);\coordinate (C) at (1.6,3);
\draw[thick,exblue] (A)--(B)--(C)--cycle;
\node[below left] at (A){#1};\node[below right] at (B){#2};\node[above] at (C){#3};
\node[below] at ($(A)!0.5!(B)$){#6};
\node[right] at ($(B)!0.5!(C)$){#4};
\node[left] at ($(A)!0.5!(C)$){#5};
\end{tikzpicture}\end{center}}

\fancyhead[R]{\footnotesize MCR3U \textperiodcentered{} 4.4}
\begin{document}

\begin{center}
{\LARGE\bfseries\color{titleblue}Worksheet 4.4 --- The Sine Law \& Cosine Law}\\[2pt]
{\small Grade 11 Functions, University (MCR3U) \textperiodcentered{} Unit 4: Trigonometric Functions}
\end{center}

\begin{ideabox}
Solve oblique triangles: sine law for a side–angle pair, cosine law for SAS or SSS.
\begin{itemize}[leftmargin=*,itemsep=2pt]
  \item Sine law: $\dfrac{a}{\sin A}=\dfrac{b}{\sin B}=\dfrac{c}{\sin C}$.
  \item Cosine law: $c^2=a^2+b^2-2ab\cos C$.
  \item SSA is the ambiguous case — it may give two triangles.
\end{itemize}
\end{ideabox}

\begin{learnbox}
\lh{Two laws for any triangle}The \textbf{sine law} $\dfrac{a}{\sin A}=\dfrac{b}{\sin B}=\dfrac{c}{\sin C}$ and the \textbf{cosine law} $a^2=b^2+c^2-2bc\cos A$ solve non-right triangles.
\lh{Which to use}Use the sine law with a matched side-angle pair (ASA, AAS, SSA); use the cosine law with SAS or SSS.
\lh{The ambiguous case}With two sides and a non-included angle (SSA) there may be zero, one, or two triangles — always check for a second possible angle.
\end{learnbox}

\sech{Worked Examples}

\begin{example}{ex0}{Example 1 --- Sine law (side)}
$A=40^\circ$, $B=60^\circ$, $a=10$. Find $b$.\soln $\dfrac{b}{\sin60^\circ}=\dfrac{10}{\sin40^\circ}\Rightarrow b=\dfrac{10\sin60^\circ}{\sin40^\circ}\approx13.5$.\gtri{$40^\circ$}{$60^\circ$}{$C$}{$10$}{$b$}{$c$}
\end{example}

\begin{example}{ex1}{Example 2 --- Sine law (angle)}
$a=8$, $b=11$, $A=35^\circ$. Find $B$.\soln $\dfrac{\sin B}{11}=\dfrac{\sin35^\circ}{8}\Rightarrow\sin B=\dfrac{11\sin35^\circ}{8}\approx0.789$. \\ $B\approx52.1^\circ$.
\end{example}

\begin{example}{ex2}{Example 3 --- Cosine law (side)}
$a=9$, $b=7$, $C=55^\circ$. Find $c$.\soln $c^2=81+49-2(9)(7)\cos55^\circ\approx130-72.3=57.7$. \\ $c\approx7.6$.
\end{example}

\begin{example}{ex3}{Example 4 --- Cosine law (angle)}
$a=5$, $b=6$, $c=7$. Find $C$.\soln $\cos C=\dfrac{25+36-49}{2(5)(6)}=\dfrac{12}{60}=0.2\Rightarrow C\approx78.5^\circ$:\gtri{$A$}{$B$}{$C$}{$5$}{$6$}{$7$}
\end{example}

\begin{example}{ex4}{Example 5 --- ASA side}
$A=50^\circ$, $B=70^\circ$, $a=12$. Find $b$.\soln $b=\dfrac{12\sin70^\circ}{\sin50^\circ}\approx14.7$.
\end{example}

\begin{example}{ex5}{Example 6 --- SAS side}
$a=8$, $b=5$, $C=60^\circ$. Find $c$.\soln $c^2=64+25-80\cos60^\circ=89-40=49\Rightarrow c=7$.
\end{example}

\begin{example}{ex6}{Example 7 --- SSS angle}
$a=4$, $b=5$, $c=6$. Find $C$.\soln $\cos C=\dfrac{16+25-36}{2(4)(5)}=\dfrac{5}{40}=0.125\Rightarrow C\approx82.8^\circ$.
\end{example}

\begin{example}{ex7}{Example 8 --- Choose the law}
Given two angles and a side, which law?\soln A side with its opposite angle is available — use the sine law.
\end{example}

\begin{example}{ex8}{Example 9 --- Choose the law}
Given three sides, which law?\soln No angle is known — use the cosine law.
\end{example}

\clearpage
\sech{Your Turn --- Show Your Work}
For each question, show all steps clearly in the space provided.

\begin{qbox}{Question 1}
$A=50^\circ$, $B=70^\circ$, $a=12$. Find $b$.\workspace{3cm}
\end{qbox}

\begin{qbox}{Question 2}
$A=45^\circ$, $B=65^\circ$, $a=14$. Find $b$.\workspace{3cm}
\end{qbox}

\begin{qbox}{Question 3}
$a=9$, $b=7$, $C=55^\circ$. Find $c$.\workspace{3cm}
\end{qbox}

\begin{qbox}{Question 4}
$a=6$, $b=8$, $c=11$. Find the largest angle.\workspace{3cm}
\end{qbox}

\begin{qbox}{Question 5}
$a=8$, $b=11$, $A=35^\circ$. Find $B$.\workspace{3cm}
\end{qbox}

\begin{qbox}{Question 6}
$a=5$, $b=6$, $C=60^\circ$. Find $c$.\workspace{3cm}
\end{qbox}

\begin{qbox}{Question 7}
$a=4$, $b=5$, $c=6$. Find $C$.\workspace{3cm}
\end{qbox}

\begin{qbox}{Question 8}
Which law: two angles and a side?\workspace{2.4cm}
\end{qbox}

\begin{qbox}{Question 9}
Which law: three sides?\workspace{2.4cm}
\end{qbox}

\begin{qbox}{Question 10}
$a=7$, $b=10$, $C=45^\circ$. Find $c$.\workspace{3cm}
\end{qbox}

\begin{qbox}{Question 11}
$A=40^\circ$, $a=10$, $b=8$. Find $B$.\workspace{3cm}
\end{qbox}

\begin{qbox}{Question 12}
$a=12$, $b=15$, $c=9$. Find the largest angle.\workspace{3cm}
\end{qbox}

\begin{qbox}{Question 13 --- Challenge}
$A=35^\circ$, $a=6$, $b=9$ (SSA). How many triangles are possible?\workspace{3cm}
\end{qbox}

\clearpage
\sech{Question Recap \& Answer Key}

\begin{recapbox}
\begin{multicols}{2}
\small
\begin{enumerate}[leftmargin=*,itemsep=3pt]
  \item $A=50^\circ$, $B=70^\circ$, $a=12$. Find $b$.
  \item $A=45^\circ$, $B=65^\circ$, $a=14$. Find $b$.
  \item $a=9$, $b=7$, $C=55^\circ$. Find $c$.
  \item $a=6$, $b=8$, $c=11$. Find the largest angle.
  \item $a=8$, $b=11$, $A=35^\circ$. Find $B$.
  \item $a=5$, $b=6$, $C=60^\circ$. Find $c$.
  \item $a=4$, $b=5$, $c=6$. Find $C$.
  \item Which law: two angles and a side?
  \item Which law: three sides?
  \item $a=7$, $b=10$, $C=45^\circ$. Find $c$.
  \item $A=40^\circ$, $a=10$, $b=8$. Find $B$.
  \item $a=12$, $b=15$, $c=9$. Find the largest angle.
  \item $A=35^\circ$, $a=6$, $b=9$ (SSA). How many triangles are possible?
\end{enumerate}
\end{multicols}
\end{recapbox}

\begin{keybox}
\begin{multicols}{2}
\small
\begin{enumerate}[leftmargin=*,itemsep=3pt]
  \item $\approx14.7$
  \item $\approx17.9$
  \item $\approx7.6$
  \item $C\approx102.6^\circ$
  \item $B\approx52.1^\circ$
  \item $\approx5.6$
  \item $C\approx82.8^\circ$
  \item sine law
  \item cosine law
  \item $\approx7.1$
  \item $B\approx31.0^\circ$
  \item $B=90^\circ$ (a $9$-$12$-$15$ right triangle)
  \item two triangles ($b\sin A\approx5.16<a<b$)
\end{enumerate}
\end{multicols}
\end{keybox}

\end{document}
